\documentclass[11pt,a4paper]{article}
\usepackage{fontspec}
\usepackage{kotex}
\usepackage{geometry}
\usepackage{graphicx}
\usepackage{xcolor}
\usepackage{tcolorbox}
\usepackage{booktabs}
\usepackage{array}
\usepackage{longtable}
\usepackage{colortbl}
\usepackage{fancyhdr}
\usepackage{titlesec}
\usepackage{hyperref}
\usepackage{emoji}
\usepackage{amsmath}
\usepackage{amssymb}
\usepackage{enumitem}

\geometry{left=25mm, right=25mm, top=30mm, bottom=30mm}

\setmainfont{Noto Sans CJK KR}
\setsansfont{Noto Sans CJK KR}
\setmonofont{Noto Sans CJK KR}
\setemojifont{Noto Color Emoji}

\definecolor{primary}{HTML}{1976D2}
\definecolor{darkbrown}{HTML}{3E2723}
\definecolor{mediumbrown}{HTML}{5D4037}
\definecolor{lightbrown}{HTML}{8D6E63}
\definecolor{cream}{HTML}{FFF3E0}
\definecolor{lightgray}{HTML}{F5F5F5}
\definecolor{tableheader}{HTML}{E8E8E8}

\titleformat{\section}
  {\Large\bfseries\color{primary}}
  {\thesection.}{0.5em}{}
\titleformat{\subsection}
  {\large\bfseries\color{darkbrown}}
  {\thesubsection}{0.5em}{}
\titleformat{\subsubsection}
  {\normalsize\bfseries\color{mediumbrown}}
  {\thesubsubsection}{0.5em}{}

\renewcommand{\contentsname}{목차}
\renewcommand{\figurename}{그림}
\renewcommand{\tablename}{표}

\tcbset{
  tipbox/.style={colback=cream,colframe=primary,boxrule=1.5pt,arc=3mm,
    left=5mm,right=5mm,top=3mm,bottom=3mm,fonttitle=\bfseries\color{primary}},
  formulabox/.style={colback=lightgray,colframe=lightgray,boxrule=0pt,arc=2mm,
    left=5mm,right=5mm,top=3mm,bottom=3mm},
  summarybox/.style={colback=white,colframe=primary,boxrule=2pt,arc=4mm,
    left=5mm,right=5mm,top=5mm,bottom=5mm}
}

\hypersetup{colorlinks=true,linkcolor=primary,urlcolor=primary,bookmarks=true,unicode=true}


\pagestyle{fancy}
\fancyhf{}
\fancyhead[R]{\small\color{lightbrown}마이크로비트 센서 검색 도구 퀵가이드 v1.0.0.0}
\fancyfoot[C]{\thepage}
\renewcommand{\headrulewidth}{0.4pt}

\begin{document}

\begin{titlepage}
\centering
\vspace*{3cm}
{\Huge\bfseries\color{primary}\emoji{magnifying-glass-tilted-left} 마이크로비트 센서 검색 도구}\\[0.5cm]
{\Huge\bfseries\color{primary} 퀵가이드}
\\[1.5cm]
{\Large 시뮬레이터 버전: v1.0.3.0}
\\[2cm]
\begin{tcolorbox}[summarybox]
\centering
{\large 이 퀵가이드는 마이크로비트 센서 검색 도구와 함께 보는 내비게이션입니다.}
\end{tcolorbox}
\vfill
{\large 사물인터넷과 빅데이터}\\[0.3cm]
{\large 청주교육대학교 컴퓨터교육과}
\vspace{1cm}
\end{titlepage}

\section*{\emoji{open-book} 이 문서 사용법}

\begin{tcolorbox}[tipbox]
\textbf{이 퀵가이드는 센서 검색 도구와 함께 보는 내비게이션입니다.}

\end{tcolorbox}

\begin{center}
\renewcommand{\arraystretch}{1.4}
\begin{tabular}{|p{9.6cm}|p{5.4cm}|}
\hline
\rowcolor{tableheader}
\textbf{상황} & \textbf{참조 섹션} \\
\hline
처음 시작 & \emoji{high-voltage} 초속성 실습 \\
\hline
검색 방법이 헷갈림 & \emoji{rocket} 빠른 시작 \\
\hline
결과를 어떻게 읽는지 모름 & \emoji{magnifying-glass-tilted-left} 관찰 포인트 \\
\hline
문제 발생 & \emoji{sos} 트러블슈팅 \\
\hline
\end{tabular}
\end{center}


\section*{\emoji{high-voltage} 1분 초속성 실습}

\begin{center}
\renewcommand{\arraystretch}{1.4}
\begin{tabular}{|p{1.5cm}|p{8.3cm}|p{5.2cm}|}
\hline
\rowcolor{tableheader}
\textbf{단계} & \textbf{행동} & \textbf{확인} \\
\hline
1 & \textbf{\emoji{speech-balloon} 텍스트로 검색} 탭에서 검색창에 \textbf{"온도"} 입력 &  \\
\hline
2 & \textbf{검색} 버튼 클릭 & 결과에 센서 목록이 나타났나요? \\
\hline
3 & "내장 센서로 충분" 또는 "Grove 센서 추천" 메시지 확인 & 추천 결과가 보이면 \emoji{party-popper} \\
\hline
\end{tabular}
\end{center}


\section*{\emoji{rocket} 빠른 시작 (2가지 방법)}

\subsection*{방법 1: \emoji{speech-balloon} 텍스트로 검색}

\textbf{"무엇을 하고 싶은지"} 키워드로 검색합니다.


1. 검색창에 키워드 입력 (예: "온도", "거리", "색깔")

2. \textbf{검색} 클릭 또는 Enter

3. 결과: 해당 키워드와 관련된 센서 목록 + 추천


\begin{tcolorbox}[tipbox, title={\emoji{light-bulb} 팁}]
\emoji{light-bulb} 추천 키워드 버튼을 활용하면 더 빠릅니다!

\end{tcolorbox}

\subsection*{방법 2: \emoji{bar-chart} 데이터로 검색}

\textbf{"어떤 데이터가 필요한지"} 체크박스로 선택합니다.


1. \textbf{\emoji{bar-chart} 데이터로 검색} 탭 클릭

2. 필요한 데이터 항목에 체크 (예: \emoji{check-mark-button} 온도, \emoji{check-mark-button} 습도)

3. 자동으로 필요한 센서 조합 추천


\begin{tcolorbox}[tipbox, title={\emoji{light-bulb} 팁}]
\emoji{light-bulb} 내장 센서로 충분한 항목은 별도로 알려줍니다. 불필요한 외장 센서를 아낄 수 있어요!

\end{tcolorbox}


\section*{\emoji{desktop-computer} 화면 구성}

\begin{center}
\includegraphics[width=0.95\textwidth]{images/S5_SensorSearch_screen.png}
\end{center}


\section*{\emoji{magnifying-glass-tilted-left} 관찰 포인트}

\begin{center}
\renewcommand{\arraystretch}{1.4}
\begin{tabular}{|p{4.2cm}|p{5.1cm}|p{5.7cm}|}
\hline
\rowcolor{tableheader}
\textbf{관찰 대상} & \textbf{찾아볼 장면} & \textbf{의미} \\
\hline
"내장 센서로 충분" 메시지 & \emoji{green-circle} 배지 표시 & 추가 센서 없이 프로젝트 가능 \\
\hline
"Grove 센서 추천" & \emoji{blue-circle} 카드 목록 & 내장으로 부족해서 외장 필요 \\
\hline
데이터 항목 수 & 체크 개수 증가 \textrightarrow{} 센서 수 증가 & "많이 측정하려면 센서도 많이 필요" \\
\hline
\end{tabular}
\end{center}


\section*{\emoji{sos} 도와주세요!}

\begin{center}
\renewcommand{\arraystretch}{1.4}
\begin{tabular}{|p{4.7cm}|p{3.6cm}|p{6.7cm}|}
\hline
\rowcolor{tableheader}
\textbf{증상} & \textbf{원인} & \textbf{해결책} \\
\hline
검색 결과가 없어요 & 키워드가 너무 구체적 & 더 일반적인 단어 사용 (예: "습기" \textrightarrow{} "습도") \\
\hline
데이터 탭에서 아무 것도 안 돼요 & 체크박스를 선택하지 않음 & 최소 1개 이상 데이터 항목 체크 \\
\hline
\end{tabular}
\end{center}


\section*{\emoji{key} 핵심 용어}

\begin{center}
\renewcommand{\arraystretch}{1.4}
\begin{tabular}{|p{2.2cm}|p{6.7cm}|p{6.0cm}|}
\hline
\rowcolor{tableheader}
\textbf{용어} & \textbf{의미} & \textbf{쉬운 설명} \\
\hline
센서 조합 & 여러 센서를 함께 사용하는 구성 & 프로젝트에 필요한 센서 세트 \\
\hline
내장 센서 & 마이크로비트 기본 탑재 & 무료, 추가 장비 불필요 \\
\hline
Grove 센서 & I2C로 추가 연결하는 외장 센서 & 더 정밀·다양한 측정 가능 \\
\hline
\end{tabular}
\end{center}


\section*{참고자료}

\begin{center}
\renewcommand{\arraystretch}{1.4}
\begin{tabular}{|p{1.5cm}|p{9.1cm}|p{2.1cm}|p{2.3cm}|}
\hline
\rowcolor{tableheader}
\textbf{\#} & \textbf{자료명} & \textbf{유형} & \textbf{비고} \\
\hline
1 & 마이크로비트센서검색도구\_v1.0.3.0 & HTML 파일 & 시뮬레이터 \\
\hline
2 & 마이크로비트센서검색도구\_개발참조\_v1.0.0.0 & 문서 & 기술 상세 \\
\hline
\end{tabular}
\end{center}


\vfill
\begin{center}
\small\color{lightbrown}
\textit{마이크로비트 센서 검색 도구 퀵가이드}
\end{center}
\end{document}